\chapter{Evaluation: Simulated scene 1} \label{ch:sim01}

The first evaluation of the atmospheric correction scheme is performed with an inland simulated scene, such that the surface and atmospheric variables are known, with a simplified spatial pattern for the surface reflectance. Calculations are performed with Eq.~\ref{eq:theory_me_closed_form} to simulate the signal that would be observed over the Blankaart reservoir in Belgium, with the Operational Land Imager (OLI) onboard Landsat 8. The scene contains only two types of targets: water and vegetation (Fig.~\ref{fig:sim01_surf_setup}A). Water is represented by the Blankaart reservoir, an hexagonal water treatment plant with a radius of 450~m, and the Blankaart lake, with a irregular shape (Fig.~\ref{fig:sim01_surf_setup}B). The distance between the water systems is $\approx$ 1~km. All other positions of the infinite plane scene represent vegetation. 

\begin{figure}[!htbp]
  \centering
    \begin{subfigure}[t]{0.48\textwidth}
      \centering
      \includegraphics[trim={0 0.75cm 0 1cm},clip,width=\linewidth]{base_spectra.pdf}
    \end{subfigure}
    \begin{subfigure}[t]{0.48\textwidth}
      \centering
      \includegraphics[trim={0 0.75cm 0 1cm},clip,width=\linewidth]{case_1_rho_s.pdf}
    \end{subfigure}
   \caption{Surface reflectance used in the simulation. (A) Spectral hemispherical-directional reflectance of surface targets, scaled by $\smash{\pi}$~sr. (B) Surface pattern of the hemispherical-directional surface reflectance scaled by $\smash{\pi}$~sr for the near-infrared waveband of OLI. In (A), subscript `w' stands for `water' and subscript `v' stands for `vegetation'. The grey polygons in the background show the relative spectral response functions of the first seven OLI wavebands.}
   \label{fig:sim01_surf_setup}
\end{figure}

The properties of the atmosphere were calculated with the 6SV code, using an atmospheric pressure at the surface of 1013.25~HPa and two standard aerosol models included in 6SV: continental and maritime. Those are the same aerosol models included in ACOLITE. The sun zenith angle was set to 32$^\circ$ and the view angle to nadir. The waveband average upward diffuse transmittance and spherical albedo for this atmosphere configuration is presented in Fig.~\ref{fig:sim01_atm_setup} for an $\smash{\tau_\text{aer}(865)}$ of 0.18~(unitless). This $\smash{\tau_\text{aer}(865)}$ corresponds to an $\smash{\tau_\text{aer}(550)}$ of 0.3~(unitless) for the continental aerosol model and to an $\smash{\tau_\text{aer}(550)}$ of 0.2~(unitless) for the maritime aerosol model.

\begin{figure}[!htbp]
  \centering
    \begin{subfigure}[t]{0.48\textwidth}
      \centering
      \includegraphics[trim={0 0.75cm 0 1cm},clip,width=\linewidth]{T_spectra_c.pdf}
    \end{subfigure}
    \begin{subfigure}[t]{0.48\textwidth}
      \centering
      \includegraphics[trim={0 0.75cm 0 1cm},clip,width=\linewidth]{T_spectra_m.pdf}
    \end{subfigure}
   \caption{Selected atmospheric properties used in the simulations. Hemispherical-directional upward diffuse transmittance of the atmosphere toward the detector scaled by $\smash{\pi}$~sr, $\smash{\pi \, T_{\text{u}_\text{dif}}(\hemiup \rightarrow \hat{\xi}_\text{D})}$, and the spherical albedo, $\smash{\rho_\text{a}(\hemiup \rightarrow \hemidn)}$, for an atmosphere with continental aerosol model (A) or maritime aerosol model (B). $\smash{\tau_\text{aer}(865)=0.18}$ for both aerosol models. The grey polygons in the background show the relative spectral response functions of the first seven OLI wavebands.}
   \label{fig:sim01_atm_setup}
\end{figure}

The integro-normalised diffuse PSF, $\smash{\hat{\curlywedge}_\text{b}^*(\hemiup \rightarrow \hat{\xi}_\text{D}; \Lambda_w; \vec{x})}$, and the integro-normalised SAF, $\smash{\hat{\curlyvee}_{\text{b}}(\hemiup \rightarrow \hemidn; \Lambda_w; \vec{x})}$, were calculated using the same standard inputs as for 6SV in terms of atmospheric vertical profile, aerosol model and concentration. The resulting hemispherical-directional environmental reflectance for the upward diffuse transmittance, $\smash{\rho_{\text{s}_{\curlywedge_\text{b}^*}}(\hemidn \rightarrow \hat{\xi}_\text{D}; \Lambda_w; \vec{x})}$, and the bi-hemispherical environmental reflectance for the spherical albedo, $\smash{\rho_{\text{s}_{\curlyvee_\text{b}}}(\hemidn \rightarrow \hemiup; \Lambda_w; \vec{x})}$, for an $\smash{\tau_\text{aer}(865)}$ of 0.18~(unitless) are presented in Fig.~\ref{fig:sim01_env_averages}.

\begin{figure}[!htbp]
  \centering
    \begin{subfigure}[t]{0.48\textwidth}
      \centering
      \includegraphics[trim={0 2cm 0 1cm},clip,width=\linewidth]{case_1_rho_env_psf_c.pdf}
    \end{subfigure}
    \begin{subfigure}[t]{0.48\textwidth}
      \centering
      \includegraphics[trim={0 2cm 0 1cm},clip,width=\linewidth]{case_1_rho_env_psf_m.pdf}
    \end{subfigure}

    \begin{subfigure}[t]{0.48\textwidth}
      \centering
      \includegraphics[trim={0 2cm 0 1cm},clip,width=\linewidth]{case_1_rho_env_sph_c.pdf}
    \end{subfigure}
    \begin{subfigure}[t]{0.48\textwidth}
      \centering
      \includegraphics[trim={0 2cm 0 1cm},clip,width=\linewidth]{case_1_rho_env_sph_m.pdf}
    \end{subfigure}

   \caption{Environmental surface reflectances for an $\smash{\tau_\text{aer}(865)}$ of 0.18~(unitless). Hemispherical-directional environmental reflectance for the upward diffuse transmittance scaled by $\smash{\pi}$~sr, that is, the weighted average surface reflectance with weights given by $\smash{\hat{\curlywedge}_\text{b}(\hemiup \rightarrow \hat{\xi}_\text{D})}$, $\smash{\pi \, \rho_{\text{s}_{\curlywedge_\text{b}^*}}(\hemidn \rightarrow \hat{\xi}_\text{D})}$, for simulations with the continental aerosol model (A) and maritime aerosol model (B). Bi-hemispherical environmental reflectance for the spherical albedo, that is, the weighted average bi-hemispherical surface reflectance with weights given by $\smash{\hat{\curlyvee}_\text{b}(\hemiup \rightarrow \hemidn)}$, $\smash{\rho_{\text{s}_{\curlyvee_\text{b}}}(\hemidn \rightarrow \hemiup)}$, for simulations with the continental aerosol model (C) and maritime aerosol model (D). All variables presented for the NIR waveband of OLI, centred at wavelength 865~nm.}
   \label{fig:sim01_env_averages}
\end{figure}

In this chapter, the near-infrared waveband of OLI, centred at 865~nm, will be used to illustrate quantities and evaluate the performance of TSDSF and RAdCor, as it is the waveband with the highest reflectance contrast between the input spectra for water and vegetation (Fig.~\ref{fig:sim01_surf_setup}), while covering an atmospheric window with significant upward diffuse transmittance for both aerosol models (Fig.~\ref{fig:sim01_atm_setup}). For this reason, specific examples use a different $\smash{\tau_\text{aer}(550)}$ for each model, such as to result in the same $\smash{\tau_\text{aer}(865)}$.

\section{Evaluation of TSDSF}

The estimation of $\smash{\rho_\text{as}^\text{Hyp}(\hat{\xi}_0 \rightarrow \hat{\xi}_\text{D}; \Lambda_w; \vec{x})}$ from $\smash{\rho_\text{as}(\hat{\xi}_0 \rightarrow \hat{\xi}_\text{D}; \Lambda_w; \vec{x})}$ with Eq.~\ref{eq:est_rho_e_toa} is an important component to find the darkest pixels at the surface and to estimate $\smash{\tau_\text{aer}(550)}$ and aerosol model. To isolate the error of the approximation described in Eq.~\ref{eq:est_rho_e_toa} from other source of errors, we let the predefined $\smash{\tau_\text{aer}(550)}$ be equal to the $\smash{\tau_\text{aer}(550)}$ used in the simulations and use a large PSF radius (25~km). Fig.~\ref{fig:sim01_hyp_estimate} shows the ratio of the estimated function, $\smash{\rho_\text{as}^\text{Hyp Est.}}$, to the reference function, $\smash{\rho_\text{as}^\text{Hyp Ref.}}$, calculated with Eq.~\ref{eq:ret_rho_env_toa} for an $\smash{\tau_\text{aer}(865)}$ of 0.18~(unitless). For those example conditions, the maximum error in the estimation was $< 10$~\% for simulations with both aerosol models. 

\begin{figure}[!htbp]
  \centering
    \begin{subfigure}[t]{0.48\textwidth}
      \centering
      \includegraphics[trim={0 2cm 0 1cm},clip,width=\linewidth]{case_1_eval_rho_env_toa_c.pdf}
    \end{subfigure}
    \begin{subfigure}[t]{0.48\textwidth}
      \centering
      \includegraphics[trim={0 2cm 0 1cm},clip,width=\linewidth]{case_1_eval_rho_env_toa_m.pdf}
    \end{subfigure}
    
   \caption{Evaluation of the estimation of $\smash{\rho_\text{as}^\text{Hyp}(\hat{\xi}_0 \rightarrow \hat{\xi}_\text{D}; \Lambda_w; \vec{x})}$ from $\smash{\rho_\text{as}(\hat{\xi}_0 \rightarrow \hat{\xi}_\text{D}; \Lambda_w; \vec{x})}$. (A) Ratio of $\smash{\rho_\text{as}^\text{Hyp Est.}}$ to $\smash{\rho_\text{as}^\text{Hyp Ref.}}$ for the simulation with the continental aerosol model; (B) same as (A) but for the simulation with the maritime aerosol model. ``Est.'' stands for estimated value and ``Ref.'' stands for reference value.}
   \label{fig:sim01_hyp_estimate}
\end{figure}

In general the error of the approximation is positive (overestimation) over the darkest pixels, due to the further blurring of the observed signal resulting from the approximation in Eq.~\ref{eq:est_rho_e_toa}. The additional blurring also helps to explain the larger error observed for the estimation of $\smash{\rho_\text{as}^\text{Hyp}(\hat{\xi}_0 \rightarrow \hat{\xi}_\text{D}; \Lambda_w; \vec{x})}$ with the maritime aerosol model simulation. The diffuse PSF for the maritime aerosol model is more peaked than the diffuse PSF for the continental aerosol model (Fig.~\ref{fig:psf_sim_res}), resulting in a smaller degree of blurring for a given $\tau_\text{aer}(550)$. Any additional blurring will have a relatively higher impact for the maritime aerosol model simulation than for the one with the continental aerosol model. The sensitivity to the additional blurring is expected to increase with the increase in $\tau_\text{aer}(550)$, in part because it decreases the relative contribution of the Rayleigh diffuse PSF, increasing the kurtosis of the mixed diffuse PSF. In addition, Eq.~\ref{eq:est_rho_e_toa} is expected to be a valid approximation in conditions where the direct upward transmittance dominates the upward transmittance. Higher values of $\tau_\text{aer}(550)$ will therefore increase the error of the approximation.

The error of the $\smash{\rho_\text{as}^\text{Hyp}(\hat{\xi}_0 \rightarrow \hat{\xi}_\text{D}; \Lambda_w; \vec{x})}$ estimation over the darkest pixels will propagate to the $\tau_\text{aer}(550)$ and aerosol model estimation with the TSDSF. First it must be considered that though the simple scenario of this simulation leaves no ambiguity of the darkest pixels at the surface, the presence of error might lead to less optimal selection of darkest surface pixels in real scenes. Further, an overestimation as seen in Fig.~\ref{fig:sim01_hyp_estimate} will lead to an underestimation $\tau_\text{aer}(550)$ for a given aerosol model. As shown in Eq.~\ref{eq:ret_ratio}, the path corrected ratio $\smash{\rho_\text{pc} / \rho_\text{pc}^\text{Hyp}}$ over the darkest pixels is equal to the diffuse fraction of the upward transmittance. Therefore, an overestimation of $\smash{\rho_\text{as}^\text{Hyp}(\hat{\xi}_0 \rightarrow \hat{\xi}_\text{D}; \Lambda_w; \vec{x})}$ over the darkest pixels will reduce the ratio, resulting in a lower diffuse fraction and a lower estimation of $\tau_\text{aer}(550)$.

We evaluate the propagated uncertainty through TSDSF using the $\tau_\text{aer}(550)$ estimation error as a metric. First we evaluate the impact of different predefined values of $\tau_\text{aer}(550)$, isolating it from errors related to the finite PSF by using a large PSF radius~(25 km). The results, presented in Fig.~\ref{fig:sim01_tsdsf_all_aot}, for a wide range of reference and predefined $\tau_\text{aer}(550)$ show two main patterns: (1) There is an overall tendency of underestimation of $\tau_\text{aer}(550)$ which is proportional to the $\tau_\text{aer}(550)$ used in the simulation. (2) The estimation improves for all simulated $\tau_\text{aer}(550)$ for larger values of predefined $\tau_\text{aer}(550)$. Further, in specific circumstances, TSDSF estimates the incorrect aerosol model for maritime aerosol simulations, resulting also in larger errors in $\tau_\text{aer}(550)$ estimation.

\begin{figure}[!htbp]
  \centering
    \begin{subfigure}[t]{0.48\textwidth}
      \centering
      \includegraphics[trim={0 0.75cm 0 1cm},clip,width=\linewidth]{case_1_eval_ts_dsf_aot_c.pdf}
    \end{subfigure}
    \begin{subfigure}[t]{0.48\textwidth}
      \centering
      \includegraphics[trim={0 0.75cm 0 1cm},clip,width=\linewidth]{case_1_eval_ts_dsf_aot_m.pdf}
    \end{subfigure}
   \caption{TSDSF estimation error of $\tau_\text{aer}(550)$. $\tau_\text{aer}(550)$ estimation error for simulations with the continental aerosol model (A) and the maritime aerosol model (B). The X-axis provides the $\tau_\text{aer}(550)$ used in the simulation, while the Y-axis provides the $\tau_\text{aer}(550)$ used for the TSDSF estimation.}
   \label{fig:sim01_tsdsf_all_aot}
\end{figure}

The errors in the estimation of $\smash{\rho_\text{as}^\text{Hyp}(\hat{\xi}_0 \rightarrow \hat{\xi}_\text{D}; \Lambda_w; \vec{x})}$ discussed previously seem to explain the observed error patterns in $\tau_\text{aer}(550)$ retrieval. The degrading performance of TSDSF with increase in the simulated $\tau_\text{aer}(550)$ is intrinsic to the TSDSF approach, as the approximation used to estimate the hypothetical quantity is expected to hold only for conditions in which the direct fraction dominates the upward transmittance. The improvement of estimation with larger predefined $\tau_\text{aer}(550)$ implies that the further blurring caused by the convolution of $\smash{\rho_\text{as}(\hat{\xi}_0 \rightarrow \hat{\xi}_\text{D}; \Lambda_w; \vec{x})}$ with the predefined PSF is minimised by reducing the relative contribution of the Rayleigh PSF. This suggests that the approximation of Eq.~\ref{eq:est_rho_e_toa} can be improved by using only the aerosol diffuse PSF, that is, without the Rayleigh contribution. The worse performance of $\tau_\text{aer}(550)$ retrievals for the simulations with the maritime aerosol model are again a consequence of the higher sensitivity to further blurring in the estimation of $\smash{\rho_\text{as}^\text{Hyp}(\hat{\xi}_0 \rightarrow \hat{\xi}_\text{D}; \Lambda_w; \vec{x})}$, as it has a higher diffuse fraction of the upward diffuse transmittance for a given $\tau_\text{aer}(550)$ value. Nonetheless, considering that the observed average $\tau_\text{aer}(550)$ is typically $< 0.3$ over most of the globe, TSDSF should provide appropriate estimates for most scenes.

The effect of the PSF extent on the TSDSF estimation was isolated by using the best predefined $\tau_\text{aer}(550)$ from the previous analysis, 1 (unitless), and evaluated by performing the TSDSF approach over a range of PSF radius from 2.5~km to 25~km. Fig.~\ref{fig:sim01_tsdsf_all_psf} shows that the $\tau_\text{aer}(550)$ estimation error is largely independent of PSF extent. The lack of sensitivity of TSDSF to the PSF radius likely results from the high coverage of the diffuse PSF when using only the PSF for aerosols (high predefined $\tau_\text{aer}(550)$). A PSF radius of 5~km covers $> 90$~\% of the diffuse transmittance for both aerosol models (Fig.~\ref{fig:psf_sim_res}). This observation can be exploited by employing a shorter PSF extent to reduce the computations for TSDSF. In addition, a shorter PSF extent resulting in lower coverage can be exploited to partially compensate for the further blurring in the estimation of $\smash{\rho_\text{as}^\text{Hyp}(\hat{\xi}_0 \rightarrow \hat{\xi}_\text{D}; \Lambda_w; \vec{x})}$, by re-normalising the truncated diffuse PSF instead of complementing the missing signal with the local average (\textit{cf.} Chapter~\ref{ch:finite}). 

\begin{figure}[!htbp]
  \centering
    \begin{subfigure}[t]{0.48\textwidth}
      \centering
      \includegraphics[trim={0 0.75cm 0 1cm},clip,width=\linewidth]{case_1_eval_ts_dsf_psf_c.pdf}
    \end{subfigure}
    \begin{subfigure}[t]{0.48\textwidth}
      \centering
      \includegraphics[trim={0 0.75cm 0 1cm},clip,width=\linewidth]{case_1_eval_ts_dsf_psf_m.pdf}
    \end{subfigure}
   \caption{TSDSF estimation error of $\tau_\text{aer}(550)$. $\tau_\text{aer}(550)$ estimation error for simulations with the continental aerosol model (A) and the maritime aerosol model (B). The X-axis provides the $\tau_\text{aer}(550)$ used in the simulation, while the Y-axis provides the PSF extent used for the TSsDSF estimation.}
   \label{fig:sim01_tsdsf_all_psf}
\end{figure}

\section{Evaluation of RAdCor}

With the estimated aerosol model and concentration, the optical functions of the atmosphere can be calculated. These, together with the estimation of $\smash{\rho_{\text{s}_{\curlyvee_\text{b}}}(\hemidn \rightarrow \hemiup; \Lambda_w; \vec{x})}$, can be plugged in to Eq.~\ref{eq:radcor_lambertian_scene_fourier_finite} to retrieve $\smash{\rho_\text{s}(\hemidn \rightarrow \hat{\xi}_\text{D}; \Lambda_W; \vec{x})}$. 

The approximations required to estimate $\smash{\rho_{\text{s}_{\curlyvee_\text{b}}}(\hemidn \rightarrow \hemiup; \Lambda_w; \vec{x})}$ are evaluated directly by setting the aerosol model and $\tau_\text{aer}(550)$ to the same values used in the simulation, and by using a PSF extent of 25~km. Fig.~\ref{fig:sim01_env_sph_est} shows an example for the simulations with $\tau_\text{aer}(550) = 0.3$ (unitless). Despite the approximations necessary to estimate $\smash{\rho_{\text{s}_{\curlyvee_\text{b}}}(\hemidn \rightarrow \hemiup; \Lambda_w; \vec{x})}$ from the TOA signal, the blurring caused by the SAF is dominant, allowing for the use of the homogeneous surface reflectance framework and resulting in a small error to the reference value. We note, however, that $\smash{\rho_{\text{s}_{\curlyvee_\text{b}}}(\hemidn \rightarrow \hemiup; \Lambda_w; \vec{x})}$ is that for the first order of the surface-atmosphere interaction.

\begin{figure}[!htbp]
  \centering
    \begin{subfigure}[t]{0.48\textwidth}
      \centering
      \includegraphics[trim={0 2cm 0 1cm},clip,width=\linewidth]{case_1_eval_rho_env_sph_c_est.pdf}
    \end{subfigure}
    \begin{subfigure}[t]{0.48\textwidth}
      \centering
      \includegraphics[trim={0 2cm 0 1cm},clip,width=\linewidth]{case_1_eval_rho_env_sph_m_est.pdf}
    \end{subfigure}
   \caption{Evaluation of the $\smash{\rho_{\text{s}_{\curlyvee_\text{b}}}(\hemidn \rightarrow \hemiup; \Lambda_w; \vec{x})}$ estimate against the reference value calculated with Eq.~\ref{eq:theory_bhem} for the OLI waveband centered at 865~nm. (A) Ratio for the Continental model simulation; (B) Ratio for the Maritime model simulation.}
   \label{fig:sim01_env_sph_est}
\end{figure}

The final critical step for the retrieval of the surface reflectance is the selection of the finite PSF extent within the constraints possible for a given scene extent (\textit{cf.} Chapter~\ref{ch:finite}). The missing fraction of the diffusely transmitted signal has to be accounted for by an estimated value, with the default in RAdCor being the local average. Therefore, the error in the correction will be proportional to the PSF extent. We isolate the effect of the PSF extent by performing the atmospheric correction using the reference $\tau_\text{aer}(550)$ and aerosol model, an evaluate the retrieval over a range of PSF radius. Fig.~\ref{fig:sim01_radcor_all} shows the average surface reflectance estimation error over the Blankaart reservoir. As expected, there is an improvement in the retrievals with larger PSF radius, with the lower boundary dependent on the $\tau_\text{aer}(550)$ used in the simulation. 

The results suggests that the largest possible PSF extent, to a maximum of 25~km, should be used for the retrieval of surface reflectance. This value, however, will be constrained by the scene size. As an example, for a centre scene area of 7~km, a PSF radius of 25~km requires a scene extent of 57~km. For targets close to the edge of the sensor swath and for commercial very high spatial resolution imagery, shorter PSF extents will be necessary, resulting in larger estimation errors.

We note, however, that the need for large kernel extents was not supported by the preliminary validation with real imagery. While other insights provided by the simulation were useful in improving the processor, kernel radius as low as 5~km or 3.5~km gave similar results to larger kernels.  

\begin{figure}[!htbp]
  \centering
    \begin{subfigure}[t]{0.48\textwidth}
      \centering
      \includegraphics[trim={0 0.75cm 0 2.2cm},clip,width=\linewidth]{case_1_eval_radcor_abs_c}
    \end{subfigure}
    \begin{subfigure}[t]{0.48\textwidth}
      \centering
      \includegraphics[trim={0 0.75cm 0 2.2cm},clip,width=\linewidth]{case_1_eval_radcor_abs_m}
    \end{subfigure}
    
    \begin{subfigure}[t]{0.48\textwidth}
      \centering
      \includegraphics[trim={0 0.75cm 0 2.2cm},clip,width=\linewidth]{case_1_eval_radcor_rel_c}
    \end{subfigure}
    \begin{subfigure}[t]{0.48\textwidth}
      \centering
      \includegraphics[trim={0 0.75cm 0 2.2cm},clip,width=\linewidth]{case_1_eval_radcor_rel_m}
    \end{subfigure}

   \caption{Evaluation of atmospheric correction over the Blankaart reservoir for different finite PSF radius. Difference between the average estimated surface reflectance and the reference reflectance for the continental model simulations (A) and maritime model simulations (B). Ratio of the average estimated surface reflectance and the reference reflectance for the continental model simulations (C) and maritime model simulations (D). The reference $\tau_\text{aer}(550)$ for each simulation was used for the correction.}
   \label{fig:sim01_radcor_all}
\end{figure}

\section{Combined evaluation of TSDSF and RAdCor}

In the previous section, RAdCor was evaluated in absence of errors in the estimation of aerosol model and $\tau_\text{aer}(550)$. Here, the aerosol properties used for the correction are taken from the TSDSF estimation. The impact of the underestimation of $\tau_\text{aer}(550)$ can be clearly seen in Fig.~\ref{fig:sim01_full_eval}. In absence of errors in $\tau_\text{aer}(550)$ estimation, the retrieved surface reflectance for the maritime aerosol model has better performance than for the continental aerosol model (Fig.~\ref{fig:sim01_radcor_all}). This is a result of the more peaked PSF for the maritime aerosol model when compared to the PSF of the continental aerosol model, causing less blurring for a given $\tau_\text{aer}(550)$. However, Fig.~\ref{fig:sim01_tsdsf_all_psf} showed the largest underestimation of $\tau_\text{aer}(550)$ for the maritime aerosol model. This also results from the shape of the PSF, as the further blurring caused by the estimation with Eq.~\ref{eq:est_rho_e_toa} leads to higher errors for a maritime model simulation. The error in the estimation of $\tau_\text{aer}(550)$ is dominant and results in lower performance for retrievals of surface reflectance for the maritime aerosol model simulations when compared to the continental aerosol model simulations. 

\begin{figure}[!htbp]
  \centering
    \begin{subfigure}[t]{0.48\textwidth}
      \centering
      \includegraphics[trim={0 0.75cm 0 1cm},clip,width=\linewidth]{case_1_full_eval_c.pdf}
    \end{subfigure}
    \begin{subfigure}[t]{0.48\textwidth}
      \centering
      \includegraphics[trim={0 0.75cm 0 1cm},clip,width=\linewidth]{case_1_full_eval_m.pdf}
    \end{subfigure}
   \caption{Evaluation of $\smash{\rho_\text{s}}$ with the combination of TSDSF and RAdCor for the OLI waveband centred at 865~nm. (A) Estimation for the continental model simulation. (B) Estimation for the maritime model simulation. The grey horizontal line provides the reference $\smash{\pi\,\rho_\text{s}(865)}$ value.}
   \label{fig:sim01_full_eval}
\end{figure}




